\documentclass[12pt,a4paper]{article}
\usepackage{mathptmx}
\usepackage[margin=2.5cm]{geometry}
\usepackage{setspace}
\onehalfspacing
\pagestyle{empty}

\begin{document}

\begin{flushright}
\today
\end{flushright}

\vspace{1cm}

\noindent Dear Professor Satchell,

\vspace{0.5cm}

Please consider the attached manuscript, ``Finite-Sample Limitations of VaR Backtesting under Ex-Ante Volatility Regimes,'' for publication in the Journal of Risk Model Validation.

Ex-ante volatility regimes for Mexican FIBRAs leave only 62--95 testable observations. At $n = 62$, the Kupiec test rejects a correctly specified model 7.6\% of the time (nominal 5\%) and has 43\% power against a model breaching at 10\%. The Christoffersen test is uninformative. The paper proposes replacing asymptotic tests with exact binomial inference and a conservatism metric $C = \hat{p} - \alpha$, neither of which depends on sample size. Across five equity tickers, GARCH-$t$ consistently calibrates near the nominal rate (pooled 3.6\%) while XGBoost-Pure over-breaches (9.7\%, $p < 0.0001$).

The manuscript has not been published elsewhere nor submitted to another journal. A separate title page with identifying information is included for the blind review process. Data and code will be released after acceptance.

The short history of many emerging-market instruments makes the finite-sample problem studied here directly relevant to risk managers and regulators working with those assets. More broadly, any validation exercise constrained by ex-ante regime definitions, new-product histories, or regulatory look-back windows faces the same structural limitation: below 100 observations, standard backtests lose discriminatory power. The decision framework proposed in Section~6 offers a practical alternative that requires no additional data and no asymptotic approximations.

\vspace{0.5cm}

\noindent Sincerely,

\vspace{1cm}

\noindent Germ\'an Pancardo

\end{document}
